\documentclass[heading=true,UTF8,a4paper,zihao=-4]{ctexart}
\usepackage{amsmath,amssymb}
\usepackage{graphicx}
\usepackage{booktabs,longtable,array}
\usepackage{xcolor}
\usepackage{geometry}
\usepackage{enumitem}
\usepackage[most]{tcolorbox}
\usepackage[colorlinks=true,linkcolor=blue,urlcolor=blue,citecolor=blue]{hyperref}
\usepackage{fancyhdr}
\usepackage{microtype}
\graphicspath{{imgs/}}

\geometry{left=2.45cm,right=2.45cm,top=2.5cm,bottom=2.4cm}
\setlist{itemsep=2pt,topsep=4pt}
\pagestyle{fancy}
\fancyhf{}
\fancyhead[L]{\small AI 智能体使用验证报告}
\fancyhead[R]{\small 校园二手交易与物品流转平台}
\fancyfoot[C]{\thepage}
\renewcommand{\headrulewidth}{0.4pt}
\setlength{\headheight}{14pt}
\setlength{\parindent}{2em}
\setlength{\parskip}{0.25em}
\renewcommand{\arraystretch}{1.18}
\definecolor{AIsectionblue}{RGB}{15,124,169}
\newcommand{\AISection}[1]{%
  \clearpage
  \refstepcounter{section}%
  \addcontentsline{toc}{section}{#1}%
  \vspace*{0.2em}\noindent
  \colorbox{AIsectionblue}{\parbox{\dimexpr\textwidth-2\fboxsep\relax}{\color{white}\bfseries\large #1}}%
  \par\vspace{0.6em}}
\newtcolorbox{QABox}{breakable,colback=white,colframe=black,boxrule=0.7pt,arc=0pt,left=8pt,right=8pt,top=8pt,bottom=8pt,before skip=0.7em,after skip=0.7em}

\usepackage{listings}
\definecolor{pygreen}{RGB}{28,172,120} % 注释颜色
\definecolor{pyblue}{RGB}{0,0,255}    % 关键字颜色
\definecolor{pyred}{RGB}{160,32,240}   % 字符串颜色

\lstset{
    language=bash,                
    basicstyle=\small\ttfamily,     
    columns=flexible,               % 3. 强烈建议改为 flexible，fixed 会导致字符间距看起来很疏离
    numbers=left,
    numberstyle=\tiny\color{gray},
    keywordstyle=\color{pyblue}\bfseries, % 4. 关键字加粗，更清晰
    commentstyle=\color{pygreen}\itshape,
    stringstyle=\color{pyred},
    showstringspaces=false,
    breaklines=true,                % 自动换行
    breakatwhitespace=false,         % 5. 只在空格处换行，避免切断变量名
    frame=shadowbox,                % 6. 换成 shadowbox 或 lines，比 single 更有质感
    rulesepcolor=\color{red!20!green!20!blue!20},
    xleftmargin=2em,                % 7. 整体左移，给行号留出空间
    framexleftmargin=2em,
    captionpos=t,                   % 标题放在上方
    aboveskip=1.5em,                % 代码块前后的间距
    belowskip=1.5em,
    escapeinside=``,                % 8. 允许在代码中插入中文（用反引号包裹）
}

% 截图放置于 imgs/ 目录；缺图时显示占位说明，避免整份报告无法编译。
\newcommand{\ReportFigure}[3]{%
\begin{figure}[htbp]
  \centering
  \IfFileExists{imgs/#1.png}{\includegraphics[width=0.82\textwidth]{imgs/#1.png}}{%
    \IfFileExists{imgs/#1.jpg}{\includegraphics[width=0.82\textwidth]{imgs/#1.jpg}}{%
      \IfFileExists{imgs/#1.pdf}{\includegraphics[width=0.82\textwidth]{imgs/#1.pdf}}{%
        \fbox{\parbox[c][3.2cm][c]{0.78\textwidth}{\centering
        \textbf{截图文件待补充}\\[0.5em]
        请将实际截图文件 \texttt{#1} 放入 \texttt{imgs/} 目录。}}}}}
  \caption{#2}
  \label{#3}
\end{figure}}
\begin{document}

\begin{center}
{\LARGE\bfseries AI 智能体使用验证报告}\\[0.65em]
{\large 项目：校园二手交易与物品流转平台}\\[0.5em]
{\normalsize 软件工程课程项目过程记录与验证说明}
\end{center}
\vspace{0.4em}
\begin{center}
\begin{tabular}{>{\bfseries}p{0.18\textwidth}p{0.28\textwidth}>{\bfseries}p{0.18\textwidth}p{0.24\textwidth}}
报告编写人 & 马歆远 & 报告日期 & 2026年10月29日\\
报告范围 & 数据与权限需求分析、管理员权限走查 & 文档性质 & 使用验证与过程记录
\end{tabular}
\end{center}
\vspace{0.5em}
\noindent\textbf{阅读说明：}本报告记录 AI 辅助需求分析、权限规则核对和走查材料整理的过程。报告中的“AI 输出摘要”用于概括相关交互，不能替代原始聊天记录；所有事实性结论仍以项目需求文档、实际截图和可复核材料为准。

\noindent\textbf{报告说明：}本报告依据本项目需求分析、管理员权限模块检查、墨刀原型走查及相关 AI 辅助讨论整理，重点说明 AI 在需求理解、权限核对、问题识别和文档整理中的实际用途，以及人工复核和责任边界。对话示例按实际工作主题归纳，用于说明人机协作过程；提交时应将实际聊天记录、原型截图和走查记录作为附件，并按证据索引补充具体文件名、日期与位置。未能通过原型或运行环境核验的功能，不表述为已实现或已验证。

\clearpage
\noindent{\Large\bfseries 目录}
\vspace{0.4em}
\tableofcontents
\clearpage

\AISection{一、AI 工具版本及具体使用情况}
\noindent\textbf{1. ChatGPT（具体模型版本以实际会话界面为准）}

\textbf{具体使用目的和环节：}ChatGPT 主要用于需求分析材料的理解与整理、管理员角色权限规则的核对、原型走查记录的规范化，以及报告结构与表述的完善。项目为“校园二手交易与物品流转平台”，本次重点工作集中在“数据与权限”相关需求，以及管理员审核工作台的原型检查。需求规格说明书 v1.1 第 5 节涉及角色判定、权限矩阵、本人和他人数据的区分及权限业务规则，因此 AI 的辅助重点并非自行创造业务规则，而是帮助把已有规则转换成可检查的问题、走查步骤和记录字段。

\textbf{主要提示方式与使用过程说明：}使用时先提供已有文档、任务分工或原型信息，再明确要求 AI 围绕指定范围分析。例如，针对管理员角色，要求逐项核对其能否浏览商品、审核商品、下架违规商品、删除他人评价、处理或驳回举报、按规则扣减信用分，同时检查管理员是否被错误赋予发布商品、下单、评价或以举报人身份举报等普通用户行为。对于原型走查，则要求区分“页面上能看到的内容”“点击后实际发生的交互”与“后台权限和数据逻辑”，避免仅凭页面外观推断系统功能已经实现。

\textbf{对 AI 输出的采纳、人工修改和核验情况：}采纳了将权限检查拆分为具体问题、以需求编号关联检查记录、把实际观察与预期行为分栏记录等建议。AI 提出的疑点仅作为候选问题，仍需回到需求规格说明书、权限矩阵、原型页面和已有走查材料进行核对。涉及管理员是否具有某项权限、操作是否会改变业务状态、是否产生操作留痕等关键结论，不能仅凭 AI 的判断确定。原型页面无法证明的后台权限和静态界面未表现的业务逻辑，均保留为“未验证”或“待确认”。

\vspace{0.4em}
\noindent\textbf{2. 墨刀（原型查看与交互走查工具）}

\textbf{具体使用目的和环节：}墨刀用于查看管理员审核工作台原型。入口为：\url{https://modao.cc/ai/share/project/6ad8788f75d0275ef61e810}。其中 \texttt{admin.html} 涉及订单审核、商品审核和举报处理。墨刀用于界面观察，不能单独证明后台权限校验、数据更新或状态转换已经实现。

\textbf{使用方式与人工核验：}按走查目标打开页面、观察操作入口并保存截图。记录中区分“页面可见”“交互可操作”和“业务逻辑已验证”。按钮或静态数据不能证明操作会正确校验权限、更新数据或生成日志；无法观察的结果标记为未验证。

\textbf{使用边界：}墨刀是原型查看工具，不等同于 AI 智能体或真实系统功能测试。最终结论由项目成员依据实际页面和证据作出。

\AISection{二、典型人机交互记录}
\noindent 以下记录按本项目实际工作主题归纳。
\ReportFigure{截图1}{任务 AI-01：需求分析任务提示词与上下文记录}{fig:ai01}

\begin{QABox}

\textbf{任务 AI-01：梳理管理员角色的权限边界}

\textbf{任务背景与输入：}
\begin{lstlisting}
任务：校园二手交易与物品流转平台——数据、权限与质量需求分析
一、你的角色
你是一名具有丰富经验的软件需求工程师、UML 建模专家和软件工程课程项目文档编写专家。

现在需要你独立完成《校园二手交易与物品流转平台》软件工程课程项目中的“数据 + 权限 + 质量”需求分析工作。

请优先阅读当前工作目录下的 马歆远-数据与权限.md 文件。这是项目组长提供的正式任务说明，也是本次工作的主要依据。

你的工作目标不是讨论应该怎么完成，而是直接执行任务，生成能够交给组长审阅的完整
初稿文件。

除非缺少的信息导致任务完全无法继续，否则不要反复询问我。可以在明确标注的前提下
使用合理假设，先完成初稿。
二、任务范围与重要限制
本次任务属于软件工程的需求分析阶段，不是系统设计或编码实现阶段。

请严格遵守以下限制：
1. 不需要开发网站或 App。
2. 不需要创建数据库或编写 SQL。
3. 不需要实现任何后端接口。
4. 不需要部署系统。
5. 不要擅自扩展为完整的软件需求规格说明书。
6. 不要修改其他组员负责的工作内容。
7. 不要引入与本项目无关的复杂技术架构。
8. 所有内容应以组长提供的 Markdown 文件为主要依据。
9. 不确定的业务规则必须明确标记为“待确认”，不要伪装成已经确定的要求。
10. 不要为了追求完整而凭空添加大量不存在的业务功能。
最终成果必须符合普通本科软件工程课程需求分析作业的要求，内容专业、逻辑清晰、格 式规范、便于验收。
三、需要完成的四项核心成果
任务一：UML 类图
根据原始 Markdown 文档中的业务说明，识别校园二手交易平台的主要实体及其关系。

至少考虑：
- 学生
- 商品
- 订单
- 分类
- 评价
- 举报记录
- 信用记录
具体实体名称以正式术语表为准。如果术语表尚不存在，暂时采用原文名称，并在说明中 标注待确认。

类图要求：
1. 采用规范的 UML 类图表达方式。
2. 每个类应包含合理的核心属性。
3. 明确类与类之间的关联关系。
4. 明确关联关系的多重性，例如 1、0..1、0..*。
5. 不要把业务流程图误画成类图。
6. 属性和关系必须与数据字典一致。
7. 不确定的关系，例如一件商品是否允许多笔历史订单，要明确标记为待确认。
8. 不要为了美观牺牲准确性。
优先使用以下方案：

方案 A：直接生成 draw.io 文件

如果具备相应工具、MCP 或文件生成能力，请生成：
- 校园二手交易平台_UML类图.drawio
- 校园二手交易平台_UML类图.png
其中 .drawio 必须能够被 diagrams.net 正常打开和编辑，PNG 应当布局整齐、中文清晰、连接线不重叠。

可以直接生成兼容 draw.io 的 XML，不必依赖网页自动化。

方案 B：生成 Mermaid 并导入 draw.io

如果无法直接生成可靠的 .drawio 文件，则生成完整的 Mermaid classDiagram 源码，保存为 .mmd 文件。

确保语法正确、中文可读，并提供将其导入 draw.io 的操作方法。

优先保证模型正确性和可编辑性。如果不能导出 PNG，请明确说明，不要假装已经生成。
任务二：数据字典
围绕系统核心实体编写完整的数据字典。

至少包括：
1. 学生
2. 商品
3. 订单
4. 分类
5. 评价
6. 举报记录
7. 信用记录
每个实体都应有独立的数据字典表格。

表格至少包含以下列：
字段名称	数据类型	业务含义	约束条件	是否必填


要求：
- 使用适合需求分析阶段的逻辑数据类型，例如文本、数值、日期时间、布尔值、枚举 、列表等。
- 不使用 VARCHAR(255)、INT(11) 这类数据库实现层面的类型。
- 编号字段应说明是否唯一、是否由系统生成。
- 枚举字段应列出可能取值。
- 合理说明字段长度、数值范围、是否允许为空等约束。
- 原文没有规定的约束值必须标记为“建议值，待确认”。
- 实体之间的关联字段必须一致。
- 学生隐私信息应与后续权限表、安全需求对应。
- 商品状态、订单状态必须与项目业务规则和状态图保持一致；尚未获取正式定义时不 得将推测状态视为已确认事实。
除了逐实体的数据字典，还应补充简短的实体关系说明。
任务三：权限矩阵
按照原始文档中的用户分类建立权限矩阵：
- 游客
- 未认证学生
- 已认证学生
- 管理员
必须覆盖原始文档列出的操作：
- 浏览商品
- 发布商品
- 下单
- 评价
- 举报
- 审核商品
- 下架违规商品
- 删除他人评价
- 查看他人信用分
如有必要，可补充与核心业务直接相关的少量操作，但必须与原有功能保持一致，并明确 区分新增建议。

权限表至少包括：
操作	游客	未认证学生	已认证学生	管理员	权限拒绝后的处理


特别要求：
1. 必须严格保留原始文件中的基础权限规定，不得擅自修改。
2. 所有拒绝操作必须给出明确提示。
3. 权限拒绝时，不得错误地改变业务数据。
4. 权限拒绝时，不得错误触发业务状态变化。
5. 应记录必要的安全审计日志。
6. 需要区分“无权操作”和“业务状态不满足条件”。
7. 对于访问本人数据和他人数据的区别，需要明确说明。
8. 对涉及隐私的信息，要遵循最小权限原则。
在权限矩阵之后，补充一小节“权限控制业务规则”，说明典型的权限拒绝与状态不变行 为。
任务四：非功能性质量需求
根据原始 Markdown 文档，编写 NFR-01 至 NFR-08 共八项质量需求。

至少涉及：
- 商品列表加载性能
- 下单响应性能
- 同时在线人数
- 个人信息保护
- 支付信息安全
- 登录安全
- 用户易用性
- 系统可用性
每条质量需求应尽可能包含：
1. 需求编号
2. 需求名称
3. 需求描述
4. 量化指标
5. 验收方式
要求：
- 保留原始文档给出的具体指标。
- NFR-03 的同时在线人数可以先暂定 500 人，但必须标记为“建议值，待确认”。
- 不允许使用“系统运行流畅”“界面友好”“安全性高”之类无法直接验证的模糊描 述。
- 对于性能要求，需要给出适当的测试前提。
- 对于安全要求，需要描述可检查的行为。
- 对于易用性要求，需要给出合理的任务完成判据。
- 对于可靠性要求，需要说明统计口径。
- 不要写 Redis、MySQL、负载均衡等实现方案。
四、跨文档一致性要求
本次任务非常重视需求之间的一致性。

请在完成四项内容后，进行一次完整检查：
1. 类图中的实体是否全部在数据字典中定义？
2. 类图中的主要属性是否与数据字典一致？
3. 数据字典的关联字段是否能支持类图中的关系？
4. 权限表中的操作是否与现有功能相符？
5. 商品和订单状态是否与已知业务规则一致？
6. 权限拒绝时，是否保证业务数据和状态不被错误修改？
7. 质量需求是否可以通过测试验证？
8. 是否存在名称不统一、字段冲突或需求编号重复的问题？
9. 是否错误引入了数据库设计、接口实现或系统开发内容？
10. 是否存在缺少依据但被写成确定事实的业务规则？
对于暂时无法解决的问题，集中记录到“待组长确认事项”中。

注意：项目组统一使用 UC、BR、ST、SC、NFR、V 编号体系。本次只为自己负责的质量需求使用 NFR-01 至 NFR-08，不得擅自重新定义其他组员的编号。
五、文件组织与最终输出
请在当前目录中创建文件夹：

数据与权限_需求分析成果/

建议生成以下文件：
数据与权限_需求分析成果/
├── 01_UML类图/
│   ├── 校园二手交易平台_UML类图.drawio
│   ├── 校园二手交易平台_UML类图.png
│   └── 类图说明.md
│
├── 02_数据字典.md
├── 03_权限矩阵.md
├── 04_非功能性质量需求.md
├── 05_待组长确认事项.md
├── 06_一致性检查报告.md
├── 数据与权限需求分析_整合初稿.md
└── README.md
如果无法生成部分图像文件，应保留准确说明，不得生成无效的空文件。

其中：

整合初稿应包含完整的四项核心成果，按以下顺序组织：
1. UML 类图及说明
2. 数据字典
3. 权限矩阵
4. 非功能性质量需求
5. 待确认问题
整合初稿应当能够直接复制到 Word 中作为正式作业使用。

各独立 Markdown 文件则方便后续修改与协作。

README 应介绍文件结构、如何打开类图、哪些内容仍需确认。
六、文档风格要求
所有文档使用简体中文。

要求：
- 学术化、规范化，但避免过度复杂。
- 符合本科软件工程课程作业风格。
- 表格排版整齐。
- 逻辑清晰，术语统一。
- 不使用夸张、口语化的表达。
- 不出现“作为一个 AI 模型”等无关内容。
- 不要编造已经执行的真实系统性能测试或安全测试。
- 区分正式要求、建议要求和待确认假设。
请确保文档内容具体，而不是只有目录和空白模板。
七、执行方式
请直接开始执行，不要仅仅向我解释应该如何完成。

具体顺序：

第一步：读取原始 Markdown 文件，提取任务要求。

第二步：检查当前目录是否有其他组员提供的术语表、用例说明、业务规则、状态图等资 料。如果存在，请阅读并统一要求；如果不存在，不要编造其内容。

第三步：梳理核心实体、字段及关联关系。

第四步：生成 UML 类图文件。

第五步：完成数据字典。

第六步：完成权限矩阵。

第七步：完成非功能性质量需求。

第八步：整合完整初稿。

第九步：执行一致性检查，修正能够确定的问题。

第十步：输出待组长确认事项和文件清单。

如果执行环境支持代码运行、Mermaid 渲染或 draw.io 文件验证，应实际检查生成文件的有效性。

对未执行的验证不要声称已经通过。

最终回复请告诉我：
1. 已完成哪些任务。
2. 生成了哪些文件。
3. 类图应该如何打开与编辑。
4. 哪些内容需要组长确认。
5. 还有哪些问题可能影响 10 月 21 日的初稿提交。
现在请读取文件并开始工作。
\end{lstlisting}

\textbf{提示内容摘要：}

\textbf{AI 输出：}
\begin{lstlisting}
已完成四项核心成果、整合初稿和一致性检查，保存在 数据与权限_需求分析成果/。
- [整合初稿](D:/Users/马歆远/Desktop/大学文件/软件工程/大作业/数据与权限_需求分析成果/数据与权限需求分析_整合初稿.md)：包含完整类图说明、数据字典、权限矩阵、NFR-01～08及待确认事项。
- 类图：[可编辑 drawio](D:/Users/马歆远/Desktop/大学文件/软件工程/大作业/数据与权限_需求分析成果/01_UML类图/校园二手交易平台_UML类图.drawio)、[高清 PNG](D:/Users/马歆远/Desktop/大学文件/软件工程/大作业/数据与权限_需求分析成果/01_UML类图/校园二手交易平台_UML类图.png)。共8个实体、52个字段、13条关联。
- 独立文件：类图说明、数据字典、权限矩阵、非功能性质量需求、待确认事项、一致性检查报告、README，另附生成脚本与校验记录，共12个文件。完整清单见 [README](D:/Users/马歆远/Desktop/大学文件/软件工程/大作业/数据与权限_需求分析成果/README.md)。
类图可在 diagrams.net／draw.io 中通过“文件 → 从以下位置打开 → 设备”打开并编辑。已完成 XML 结构、关联字段、36格基础权限、编号和 PNG 视觉检查；尚未在 draw.io 编辑器内实际打开验证，也未执行真实系统测试。

需组长优先确认：正式术语、订单状态、同商品历史及有效订单限制、分类反向多重性、评价规则、举报范围、信用规则，以及500人容量和验收口径。详见[待确认事项](D:/Users/马歆远/Desktop/大学文件/软件工程/大作业/数据与权限_需求分析成果/05_待组长确认事项.md)。

影响10月21日提交的主要问题是术语表和状态资料尚未提供。建议10月18日会议后优先对齐；未解决项已明确标注，不影响本稿先交组长审阅。另需保存本次实际提示词及对话截图或链接，供后续AI使用验证报告使用。
\end{lstlisting}

\textbf{人工核验与结论：}本节应结合实际文件和原始对话补充人工核验结果；当前未提供独立的人工复核记录时，不将 AI 输出本身视为核验依据。
\end{QABox}

\ReportFigure{截图2}{任务 AI-02：提交前质量验收提示词与输出记录}{fig:ai02}

\begin{QABox}
\textbf{任务 AI-02：区分“页面存在”与“权限功能已实现”}

\textbf{任务背景与输入：}
\begin{lstlisting}
现在项目文件已经全部生成，请停止新增功能，转入严格的质量验收阶段。

请以原始文件《马歆远-数据与权限.md》为最高依据，对刚才生成的所有成果进行一次完整审查。

重点检查：
1. 文件完整性：所有要求的文件是否存在，内容是否完整。
2. UML 类图：实体是否齐全，属性是否合理，关联关系和多重性是否正确，.drawio 是否可以正常解析，PNG 是否已生成。
3. 数据字典：是否遗漏字段，字段类型和约束是否合理，实体名称是否统一。
4. 权限矩阵：是否完整覆盖原文中的四类用户及九项操作，权限拒绝后是否有明确提示、状态不变和日志记录。
5. 质量需求：NFR-01 至 NFR-08 是否完整，数值与原始要求是否一致，指标是否具有可测试性。
6. 一致性：类图、数据字典、权限表与质量需求是否存在相互矛盾的地方。
7. 需求越界：是否擅自加入组长没有确定的业务功能、状态或规则。
8. 编号规范：是否使用项目组规定的术语与编号体系。
请对发现的问题按以下标准分类：
- 严重问题：会导致作业要求不满足或明显逻辑错误。
- 一般问题：内容不一致、字段或约束不清晰。
- 建议优化：不影响提交，但可以进一步完善。
- 待组长确认：缺少外部业务定义，不能自行决定。
对于能根据原始文件明确修复的问题，请直接修复；对于缺少依据的问题，只做标注，不得自行编造。

最终生成一份《提交前质量验收报告.md》，包括问题位置、问题描述、严重程度、处理结果和待确认事项。

最后明确给出：当前成果是否达到“可以提交组长审阅的初稿”标准。

注意：不能仅根据已有的一致性检查报告宣称合格，必须重新检查实际文件。不得声称完成了没有实际执行的验证。
\end{lstlisting}
\textbf{提示内容摘要：}

\textbf{AI 输出：}
\begin{lstlisting}
已重新检查实际文件，生成 **[提交前质量验收报告.md](D:/Users/马歆远/Desktop/大学文件/软件工程/大作业/数据与权限_需求分析成果/提交前质量验收报告.md)**。

**结论：当前成果达到“可以提交组长审阅的初稿”标准。**

- 严重问题：0项。
- 一般问题：4项，已修订并同步整合稿，涉及列表加载范围、举报自动下架表述、管理员隐私访问边界及拒绝发布不得新增商品。
- 修订后82项文件与一致性检查全部通过；原文36格权限和八项NFR的已定指标保持一致。
- 正式术语、订单状态、分类反向多重性、评价与举报规则、部分验收口径仍待组长确认，未自行改定。

`.drawio` 已实际解析，PNG已重新查看。尚未执行draw.io编辑器内打开编辑及真实系统测试，报告中已明确注明。
\end{lstlisting}
\textbf{人工核验与结论：}本节应结合实际文件和原始对话补充人工核验结果；当前未提供独立的人工复核记录时，不将 AI 输出本身视为核验依据。
\end{QABox}

\ReportFigure{截图3}{任务 AI-02：质量验收结论的对话记录}{fig:ai02-result}

\begin{QABox}
\textbf{任务 AI-03：权限模块走查与交付材料整理}

\textbf{任务背景与输入：}
\begin{lstlisting}
任务：完成权限模块走查与交付材料整理
你现在需要协助我完成项目中的权限模块走查任务。请先阅读项目中的 修改要求.md，严格按照其中的要求执行。
一、总体要求
1. 先检查当前工作区的文件结构、现有需求文档、数据字典、权限矩阵、原型 HTML 文件及已有走查记录。
2. 以 修改要求.md 为本次任务的主要依据，保留其中的术语、编号、判定标准和交付要求。
3. 尽量在现有文件上修改，不要无故新建重复文件，不要大规模重构项目，不要覆盖与本次任务无关的内容。
4. 先检查现状，再决定需要修改什么。对于已有且符合要求的内容，不要重复制作。
5. 不得伪造实际测试结果、截图、运行状态或验证证据。无法通过静态分析确认的内容必须明确标注，不能猜测。
6. 所有修改应尽量与现有需求规格说明书、数据字典、权限矩阵及组长已经确定的需求保持一致。
二、检查并完成 SC-02 权限流程走查
以管理员视角检查现有原型与需求，重点核对以下事项：
1. 商品审核：只有管理员可以审核商品；管理员不能发布商品、下单或评价。
2. 举报处理：管理员不能通过普通用户举报入口举报，应直接进行管理处置。
3. 强制下架：适用对象应为在售或已预订商品；待审核商品应通过审核拒绝处理。
4. 扣信用分：应作为独立处置手段，与驳回举报、强制下架并列。
5. 撤销处置：错误处置应能够恢复，并保留操作记录。结合 BR-23、ST-28 检查对应需求和原型是否一致。
请查找项目中实际存在的页面、按钮、状态及相关业务逻辑，记录每一项的依据。

如果原型缺少对应流程，应记录缺失情况，不要直接宣称业务功能已经实现。
三、完成权限六问核对表
逐条检查下列权限规则，并结合现有原型、需求文档和代码填写核对结果：
1. 游客不能发布商品。
2. 未认证学生不能下单，应该提示“请先完成认证”。
3. 已认证学生可以查看他人信用分数，但不能查看他人的信用分变动明细。
4. 管理员不能举报，应直接进行管理处置。
5. 卖家只能下架自己发布的商品，不能下架他人商品。
6. 管理员不能下单、发布商品或评价。
核对表至少应包含：
- 权限问题
- 需求规定
- 原型实际情况
- 是否一致
- 证据或备注
注意区分以下两类拒绝：
- 角色无权：当前用户角色本身不允许执行该操作。
- 状态不允许：当前角色可能具有操作权限，但当前业务状态不允许执行。
不得将两者混为一谈。
四、完成 PV-02 走查记录
按照 修改要求.md 中的原始表格，逐项处理 PV-02-1 至 PV-02-6。

必须保留原有的验证编号、测试条件与输入、预期结果，不得擅自修改判定标准。

重点检查：
- PV-02-1：admin.html 是否只有“订单审核”“商品审核”两个 Tab；若存在“用户认证审核”，应按任务要求判定。
- PV-02-2：查找举报处理入口。
- PV-02-3：查找强制下架、驳回举报、扣信用分按钮。
- PV-02-4：核对后台登录身份显示，检查是否错误显示为学生身份。
- PV-02-5：检查首页和详情页中，未登录或未认证时按钮的置灰、提示等页面表现。
- PV-02-6：检查详情页的举报入口是否可操作。
填写时严格遵循：
1. 实际结果只能使用“通过”“失败”“未验证”三种结论，并附一句具体说明。
2. “发现的问题”没有问题时填写“无”；没有完成验证时填写“未验证”，不得混用。
3. “证据位置”按照原任务要求保留为空，供指定人员后续补充截图。
4. 页面上存在置灰或提示，只能用于判断页面层表现，不能证明后端权限拦截有效。
5. 如果只是静态 HTML 原型，没有实际登录态或权限控制机制，那么真实业务权限拦截应判为“未验证”，而不是直接判定为“失败”。
6. 如果代码中确实存在可执行的权限逻辑，可以进一步静态分析或运行验证，但必须说明验证依据和范围。
如果当前环境允许实际打开或运行原型，请尽可能完成真实页面检查。无法运行的部分应列出原因及后续人工验证步骤。
五、整理权限相关问题清单
结合上述走查结果，整理独立的问题清单。

每个问题尽量包含：
- 问题编号
- 问题描述
- 对应需求或规则
- 涉及页面、文件或功能
- 问题类型（角色无权、状态不允许、页面表现不符、业务层未验证等）
- 严重程度（有依据时再填写）
- 建议处理方式
- 验证状态
不得把尚未验证的事项直接认定为缺陷，也不要为了让表格完整而编造问题。
六、检查并完善 A.8 个人贡献记录
在 修改要求.md 中找到“A.8”对应内容，并检查其与项目实际文件、修改记录和本次交付物是否一致。

原有记录涉及：
- UML 类图、9 个实体的数据字典、权限矩阵、NFR-01～NFR-08。
- 需求规格说明书 v1.0 的对应章节。
- 合稿检查中发现的信用分公开范围、举报主体等需求矛盾。
- 订单状态枚举及相关字段、实体补充。
- .md 文件作为唯一准据、废弃脚本生成方式的调整。
- 本次 SC-02、SC-06、SC-10、PV-02 的协作与证据编号。
尽可能保留已有的事实性记录。只有在核实项目文件后，才修改与实际不符的描述。不要声称完成了尚未完成的走查或不存在的产物。
七、修改策略
请按以下顺序执行：

第一步：扫描项目，找出相关文件，阅读 修改要求.md 以及相关需求、原型和数据文件。

第二步：生成简短的检查清单，明确哪些内容已经具备、哪些需要补充、哪些需要实际运行验证。

第三步：检查原型页面和现有实现。优先通过代码、页面结构及可运行的测试确认事实，不要只根据文件名推断。

第四步：直接修改需要交付的现有文件。表格结构尽量沿用原格式；如需新建问题清单或权限核对表，优先放在项目现有的相关文档目录中。

第五步：检查所有修改后的文件，确认编号一致、术语统一、没有遗漏交付物、没有虚构验证结果，也没有覆盖无关内容。

第六步：向我汇报完成情况，包括：
- 修改了哪些文件及其路径。
- 每个文件具体改了什么。
- 六项 PV-02 检查各自的结论。
- 权限六问的核对结果。
- 当前发现的问题。
- 哪些事项仍然需要人工操作或补充截图。
- 是否有未完成的交付物，以及原因。
八、最终交付要求
目标交付物为：
1. 权限一致性核对表。
2. PV-02 走查记录。
3. 权限相关问题清单。
4. 更新后的 A.8 个人贡献记录。
优先在现有文件中完成以上交付物。如果当前目录没有适合存放某份材料的位置，再创建必要的新文件，并在最终汇报中注明路径。

现在开始检查工作区并执行任务。不要只给我修改建议，要在确认文件和环境允许的前提下，实际完成能完成的修改。遇到无法验证的内容，明确记录，不要编造结果。
\end{lstlisting}

\textbf{提示内容摘要：}
\textbf{AI 输出及人工处理：}
\begin{lstlisting}
已完成文档走查与可交付材料整理。PV-02 的证据位置按要求留空；由于工作区没有原型页面或可运行的业务代码，页面与真实权限拦截都不能据此判定。

## 修改文件

- [修改要求.md](<D:/Users/马歆远/Desktop/大学文件/软件工程/大作业/数据与权限_需求分析成果 第二次修改版/修改要求.md>)：填写 PV-02-1～PV-02-6 的实际结果和问题栏；保留原验证编号、测试条件、预期结果及空白证据位置；核实并更新 A.8。
- [权限模块走查与交付材料.md](<D:/Users/马歆远/Desktop/大学文件/软件工程/大作业/数据与权限_需求分析成果 第二次修改版/权限模块走查与交付材料.md>)：新增 SC-02 依据、权限六问核对表、问题清单及后续人工验证步骤。

静态检查确认 PV-02 六行均为六列，实际结果只使用“未验证”，发现的问题栏按规则填写“未验证”，证据位置均为空；权限六问六项、问题清单五项均已写入。

## PV-02 结论

| 编号 | 实际结果 |
|---|---|
| PV-02-1 | **未验证**：没有 `admin.html`，无法检查后台 Tab。 |
| PV-02-2 | **未验证**：没有后台原型，无法查找举报处理入口。 |
| PV-02-3 | **未验证**：没有后台原型，无法检查三种处置按钮。 |
| PV-02-4 | **未验证**：没有后台页面，无法检查管理员身份显示。 |
| PV-02-5 | **未验证**：没有首页／详情 HTML，页面置灰或提示未检查；真实权限拦截也未验证。 |
| PV-02-6 | **未验证**：没有详情页，无法检查举报入口及其可操作性。 |

## 权限六问

需求文档支持：游客不能发布；未认证学生不能下单；已认证学生能查看他人分数但不能查看变动明细；卖家只能下架自己的商品；管理员不能下单、发布或评价。以上项目的原型表现和业务层拦截均未验证。

管理员不能举报这一规则在权限矩阵和数据字典正文中有依据，但 UML 说明和 `.drawio` 仍保留管理员作为举报人的关系，存在跨文件冲突。另有两处规则细节待跟进：强制下架的商品状态条件尚未明确写入现有规则；未认证下单提示在矩阵中为“请先完成学生认证”，与本任务要求的“请先完成认证”文案不完全一致。

## 问题与未完成事项

问题清单记录了 UML 举报主体冲突、强制下架状态边界缺失、处置撤销与留痕未明确、原型材料缺失，以及提示文案差异。原型材料缺失只表示无法验证，不代表功能缺陷。BR-23、ST-28 的正文也未在当前工作区找到，因此撤销恢复和操作记录无法核实。

文档交付物已完成；实际页面走查和截图仍需人工补做：提供原型或运行地址及适当测试身份后，分别检查页面表现与业务层权限，再由指定人员补齐 PV-02 的证据位置。
\end{lstlisting}

\textbf{人工核验与结论：}本节应结合实际文件和原始对话补充人工核验结果；当前未提供独立的人工复核记录时，不将 AI 输出本身视为核验依据。
\end{QABox}

\ReportFigure{截图4}{任务 AI-03：权限模块走查任务提示词记录}{fig:ai03}

\begin{QABox}
\textbf{任务 AI-04：依据组长意见完成最终整理}

\textbf{任务背景与输入：}
\begin{lstlisting}
组长已经确认：找不到原型页面的走查项目就填写“未验证”，不需要我自行创建 admin.html 或其他新原型。

请基于当前工作区继续完成最终整理：
1. 保留 PV-02-1 至 PV-02-6 的“未验证”结论，并在说明中明确原因是当前工作区及已检查的仓库目录中没有找到对应原型页面。
2. 不要将“未验证”改成“失败”或“通过”，也不要编造页面表现和测试证据。
3. 检查并完善权限一致性核对表、PV-02 记录、权限相关问题清单及 A.8 个人贡献记录。
4. 核对管理员举报权限、强制下架商品状态范围、撤销处置与操作留痕、未认证提示文案这四项问题。能根据现有需求依据确定的就修订；需要组长决定的就保留为待确认事项。
5. 检查上一轮修改的文件是否放在正确目录，避免重复文件、互相矛盾的结论或覆盖无关内容。
6. 直接修改需要修改的文件。完成后汇报修改文件清单、主要修改内容、未解决事项和最终交付物。
请以我负责的“数据与权限_需求分析成果 第二次修改版”为工作范围，不需要整个小组的原型代码，也不要自行开发新页面。
\end{lstlisting}

\textbf{提示内容摘要：}
\textbf{AI 输出及人工处理：}
\begin{lstlisting}
收到。我会按组长的确认把 PV-02 的“未验证”保留为最终结论，并在原因中明确写明已检查的工作区／项目目录没有原型页面。现在继续修订能由现有规则确定的文档冲突：管理员举报权限、强制下架状态边界和下单提示文案；处置撤销按已确认的“可恢复并留痕”补入规则，同时把仍缺少的具体规则留作待确认。




Ran $root='D:\Users\马歆远\Desktop\大学文件\软件工程\大作业\数据与权限_需求分析成果 第二次修改版'; foreach($f in @('03_权限矩阵.md','02_数据字典.md','01_UML类图\类图说明.md','数据与权限需求分析_整合初稿.md')){ '--- '+$f+' ---'; $p=Join-Path $root $f; $n=0; Get-Content -LiteralPath $p -Encoding UTF8 | ForEach-Object {$n++; if(($f -eq '03_权限矩阵.md' -and $n -le 100) -or ($f -eq '02_数据字典.md' -and ($n -le 15 -or ($n -ge 38 -and $n -le 45) -or ($n -ge 75 -and $n -le 102) -or ($n -ge 124 -and $n -le 150))) -or ($f -eq '01_UML类图\类图说明.md' -and ($n -le 7 -or ($n -ge 30 -and $n -le 43))) -or ($f -eq '数据与权限需求分析_整合初稿.md' -and ($n -le 55 -or ($n -ge 230 -and $n -le 315)))){ '{0,4}: {1}' -f $n,$_ }} }; '--- render tools ---'; Get-Command draw.io,drawio,diagrams.net -ErrorAction SilentlyContinue | Select-Object Name,Source; '--- project root metadata ---'; Get-Item -LiteralPath $root | Select-Object FullName,Attributes; 'git directory exists: ' + (Test-Path -LiteralPath (Join-Path $root '.git'))




Viewed an image



Ran commands



我在同步 UML 和整合稿时发现 PowerShell 把我临时使用的 R 名称解释成了历史命令别名，导致那一批替换没有应用；权限矩阵和数据字典的修改已成功。我会改用无冲突的替换函数重做这两份文件，并逐项检查命中数后再继续。




Ran commands



Viewed an image



Ran commands



Viewed an image



Ran commands
\end{lstlisting}

\textbf{人工核验与结论：}本节应结合实际文件和原始对话补充人工核验结果；当前未提供独立的人工复核记录时，不将 AI 输出本身视为核验依据。
\end{QABox}

\ReportFigure{截图5}{任务 AI-04：依据组长意见继续整理交付材料的记录}{fig:ai04}

\begin{QABox}
\textbf{任务 AI-05：最终交付汇总与验证边界说明}

\textbf{任务背景与输入：}
\begin{lstlisting}
继续上面的任务
\end{lstlisting}

\textbf{提示内容摘要：}
\textbf{AI 输出及人工处理：}
\begin{lstlisting}
完成情况
已在“数据与权限_需求分析成果 第二次修改版”目录内完成文档整理，没有新建原型页面。新增的总交付材料是 [权限模块走查与交付材料.md](D:/Users/马歆远/Desktop/大学文件/软件工程/大作业/数据与权限_需求分析成果 第二次修改版/权限模块走查与交付材料.md)；其中包含权限六问核对表、SC-02 走查、问题清单和后续人工验证事项。

修改文件如下：
- [修改要求.md](D:/Users/马歆远/Desktop/大学文件/软件工程/大作业/数据与权限_需求分析成果 第二次修改版/修改要求.md)：保留 PV-02 原编号、条件和预期结果；六项实际结果均为“未验证”，问题栏均为“未验证”，证据位置留空；更新 A.8。
- [03_权限矩阵.md](D:/Users/马歆远/Desktop/大学文件/软件工程/大作业/数据与权限_需求分析成果 第二次修改版/03_权限矩阵.md)、[02_数据字典.md](D:/Users/马歆远/Desktop/大学文件/软件工程/大作业/数据与权限_需求分析成果 第二次修改版/02_数据字典.md)、[数据与权限需求分析_整合初稿.md](D:/Users/马歆远/Desktop/大学文件/软件工程/大作业/数据与权限_需求分析成果 第二次修改版/数据与权限需求分析_整合初稿.md)：同步管理员举报权限、强制下架状态、处置撤销留痕及未认证下单提示规则。
- [类图说明.md](D:/Users/马歆远/Desktop/大学文件/软件工程/大作业/数据与权限_需求分析成果 第二次修改版/01_UML类图/类图说明.md)、[校园二手交易平台_UML类图.drawio](D:/Users/马歆远/Desktop/大学文件/软件工程/大作业/数据与权限_需求分析成果 第二次修改版/01_UML类图/校园二手交易平台_UML类图.drawio)、[校园二手交易平台_UML类图.png](D:/Users/马歆远/Desktop/大学文件/软件工程/大作业/数据与权限_需求分析成果 第二次修改版/01_UML类图/校园二手交易平台_UML类图.png)：移除管理员作为举报人的关系，保留学生提交、管理员处理的关系，并同步配图。
目录递归检查未找到 HTML 原型或可运行的业务前后端代码。因此，以下 PV-02 项目均保持“未验证”：
编号	结论	原因
PV-02-1	未验证	未找到 admin.html，无法检查后台 Tab。
PV-02-2	未验证	未找到后台原型页面，无法查找举报处理入口。
PV-02-3	未验证	未找到后台原型页面，无法检查处置按钮。
PV-02-4	未验证	未找到后台原型页面，无法检查身份显示。
PV-02-5	未验证	未找到首页／详情原型或可运行应用；页面表现和真实权限拦截均未验证。
PV-02-6	未验证	未找到详情页原型，无法检查举报入口及其可操作性。


权限核对结果
六问的需求规则已纳入核对表。游客发布商品、管理员发布／下单／评价、管理员通过举报入口举报、卖家下架他人商品均按角色或对象权限拒绝；未认证学生下单按认证状态拒绝，提示文案已统一为 “请先完成认证”。已认证学生可查看他人信用分、不能查看变动明细的规则也已记录。

这些是文档核对结果，不代表页面或业务权限已实测。另有一项范围待确认：权限矩阵允许校内任意学生查看他人分数，而本次规则明确了已认证学生可查看，但没有明确排除未认证学生；本轮未擅自收窄矩阵范围。
待确认与验证边界
- 强制下架规则已明确为只适用于“在售”或“已预订”商品；“待审核”走审核拒绝，其他状态不直接强制下架。已预订商品关联订单的处理影响仍待组长／业务规则负责人确认。
- 错误处置的最低规则已补入文档：恢复其造成的业务影响、保留原处置记录并追加撤销记录。多次处置及后续状态冲突、审计关联字段和保留期限待确认；当前目录未找到 BR-23、ST-28 正文。
- 实际页面、按钮、提示和权限拦截均未验证；PV-02 证据栏仍为空，待原型／运行环境及指定截图人员补充。
- UML 源图 XML 静态解析通过；未在 diagrams.net 中打开或导出验证。未运行业务系统测试，也未制作或声称存在截图。
文档静态核对确认六条 PV 记录、空证据栏和走查报告表格结构符合要求。此前的 05、06 及提交前验收记录有历史版本说明，本轮未覆盖这些历史材料。
\end{lstlisting}

\textbf{人工核验与结论：}本节应结合实际文件和原始对话补充人工核验结果；当前未提供独立的人工复核记录时，不将 AI 输出本身视为核验依据。
\end{QABox}



\AISection{三、智能体任务约束与执行记录}
\noindent\textbf{记录类型：}实际使用中的任务约束整理。若课程将 B5 定义为必须执行的独立智能体任务，应结合教师要求确认本节是否满足其具体范围；本报告不把普通问答式辅助虚构为自主执行的智能体。

\textbf{任务目标：}在不改变既有业务规则的前提下，辅助检查校园二手交易与物品流转平台的管理员权限需求和原型走查材料，形成可复核的问题清单与报告记录。

\textbf{输入材料与阅读要求：}
\begin{enumerate}
\item 以需求规格说明书 v1.1 第 5 节及对应权限矩阵为主要业务依据。
\item 参考 PV-02 管理员权限原型走查记录、墨刀管理员工作台页面及实际截图。
\item 使用材料时注明版本和位置；无法确认版本或来源的内容不得作为确定规则。
\end{enumerate}

\textbf{允许操作范围：}可以归纳现有要求、提出核查问题、建议表格字段、整理提示词与输出摘要、检查术语和编号一致性，并对材料中缺少依据的结论提出警示。

\textbf{禁止操作范围：}不得自行新增或更改业务权限；不得将 AI 推断写成业务方确认；不得将静态页面写成真实功能已实现；不得把未执行的测试写成通过；不得虚构对话原文、截图、日志、人员确认或复验结果；不得擅自把待确认事项改写为已确认规则。

\textbf{输出形式及业务规则：}输出应包括任务目的、输入依据、关键建议、人工核验、最终决定、证据位置和遗留问题。管理员可执行的操作及禁止操作应以需求文档为准；当规则存在缺失或冲突时，应标注待确认并返回团队讨论。

\textbf{人工确认责任：}AI 只提供整理和检查建议。需求规则是否成立、问题是否修改、修改是否符合项目范围，由小组成员负责；涉及真实业务方或教师确认的内容，只有在实际确认发生后才能如此记录。团队内部复核应明确标为团队内部复核。

\textbf{信息不足时的处理：}遇到文档缺失、规则编号无法对应、原型无法显示后台结果或实际行为无法复现时，记录“待确认”或“未验证”，注明需要补充的材料及责任人，不根据一般经验补全项目事实。

\textbf{完成标准与验收方法：}
\begin{enumerate}
\item 管理员允许与禁止的操作均能在需求文档或权限矩阵中找到对应依据，或被明确标注为待确认。
\item 每条走查记录包含具体检查对象、预期行为、实际观察和结论。
\item 页面展示、交互行为与后台业务验证分别描述，不互相替代。
\item 报告中引用的证据编号能够定位到真实文件、截图、对话或需求段落。
\item 对尚未执行的测试、未确认的规则和未完成的复核保留真实状态。
\end{enumerate}

\textbf{实际执行与检查结果：}本次工作中，AI 被用于权限清单整理、走查记录结构化、风险提示和报告文字组织；项目成员通过阅读需求材料、查看墨刀原型和核对截图完成必要的人工检查。管理员页面的部分界面信息能够通过原型观察，但后台权限拦截、业务数据持久化、信用分变化和操作日志等行为，若无可执行系统或相应日志证据，均不能据此认定通过。任务约束在实际使用中主要体现为对范围和结论措辞的限制；最终提交前仍需检查附件是否完整、编号是否一致。

\textbf{必要的约束改进：}后续使用时，应在提示词中更明确地要求 AI 对每项结论标注依据位置和证据类型，并将输出分为“文档已支持”“AI 提出的候选意见”“待人工确认”三类。对测试类任务，还应要求明确区分“未执行”“执行但失败”“执行通过”，避免将没有观察到问题误写为通过。

\AISection{四、关键结论的性质与核验}
\begin{longtable}{|p{0.29\textwidth}|p{0.15\textwidth}|p{0.25\textwidth}|p{0.22\textwidth}|}
\hline
\textbf{关键结论} & \textbf{性质} & \textbf{依据} & \textbf{核验状态及结果}\\
\hline
管理员具有商品浏览、商品审核、违规商品下架等管理权限 & 需求事实 & 需求规格说明书 v1.1 第 5 节及权限矩阵 & 需与提交版本逐项核对；以文档原文为准\\
\hline
管理员不能发布商品、下单、评价或以举报人身份举报 & 权限边界 & 需求规格说明书 v1.1 第 5 节及角色规则 & 需确认相关规则在最终版本中表述一致\\
\hline
墨刀中存在管理员审核工作台页面 & 原型观察 & 墨刀原型及截图 & 页面可见性可由截图支持，不代表后台功能通过\\
\hline
后台权限校验和数据状态变更正确 & 功能结论 & 需要可执行系统、日志或有效测试证据 & 证据不足时标记“未验证”\\
\hline
AI 提出的规则冲突或遗漏 & 候选问题/建议 & AI 交互记录及对应需求材料 & 必须由人工返回原文核查，不直接视为事实\\
\hline
\end{longtable}

\AISection{五、证据索引}
\noindent 下表为建议采用的编号结构。提交前应将“待补充”替换为真实附件名称、链接和具体位置；不存在或未保存的材料不得虚构。

\begin{longtable}{|p{0.12\textwidth}|p{0.31\textwidth}|p{0.32\textwidth}|p{0.17\textwidth}|}
\hline
\textbf{证据编号} & \textbf{类型及内容} & \textbf{文件或链接及具体位置} & \textbf{支持的任务}\\
\hline
E-01 & AI 关键提示词与输出片段 & 实际 ChatGPT 对话记录；补充对话日期、轮次或导出文件 & AI-01、AI-05\\
\hline
E-02 & 需求依据与权限矩阵 & 需求规格说明书 v1.1，第 5 节及对应权限规则；补充最终文件路径 & AI-01、AI-04\\
\hline
E-03 & 管理员原型页面 & 墨刀原型链接及 \texttt{admin.html} 相关页面截图 & AI-02、AI-03\\
\hline
E-04 & PV-02 原型走查截图与记录 & 实际走查记录 Markdown 及截图文件；补充具体文件名、编号和位置 & AI-02、AI-03\\
\hline
E-05 & 问题清单、修改对照或复核记录 & 对应需求分析成果、问题清单及最终修订版本；未完成项注明待复核 & AI-03、AI-04、AI-05\\
\hline
\end{longtable}

\AISection{六、总结与责任说明}
本次 AI 辅助工作的价值在于把管理员权限拆解为可检查的操作和数据边界，将预期行为、实际观察与结论分开记录，并通过编号关联需求、走查和证据。AI 可以协助发现潜在遗漏、整理材料和改善表达，但不能代替需求依据、实际测试或团队决策。

原型页面存在不代表后台权限已实现，按钮能够点击也不代表业务数据正确更新，AI 提出问题更不代表问题已被证实。提交前应核对需求文档最终版本、PV-02 走查记录、截图和 AI 对话附件；对无法验证的项目保留“未验证”状态。

\noindent\textbf{提交前检查清单：}
\begin{enumerate}
\item 补齐小组编号、成员名单、报告日期和需求文档最终版本位置。
\item 将实际 AI 模型或版本按会话界面填写；不清楚时填写“未记录”。
\item 将 AI-01 至 AI-05 的摘要与实际聊天记录核对，必要时调整为准确的实际过程。
\item 确保证据索引中的文件真实存在、能够打开，并与报告版本一致。
\item 检查所有“通过”“已修改”“已确认”等措辞是否有实际证据；没有证据的改为“未验证”“待复核”或“待确认”。
\end{enumerate}

\end{document}
